% GENERATED FILE. Edit canonical lessons, then run npm run generate:books.
% canonical-source-hash: fnv1a64:90780c5d4b5599a0
% canonical-lessons: AR-C490-dahsha, AR-C490-mundahish, AR-C490-indahasha, AR-C490-tafajaa, AR-C490-ijab

\chapter{Surprise and admiration}
\label{ch:ar-surprise-and-admiration}
\hlchaptermodality{490}

\begin{chapteropening}
\textbf{By the end of this chapter:} \emph{I can talk about astonishment and admiration, and say that someone was amazed or surprised.}

Complete the last lesson of chapter 490: I can talk about astonishment and admiration, and say that someone was amazed or surprised.
\end{chapteropening}

\section[\texorpdfstring{\ar{دهشة} (dahsha)}{دهشة (dahsha)}]{\ar{دهشة} (dahsha) — astonishment}

\label{lesson:AR-C490-dahsha}



\begin{quote}
Before the new one: say the Arabic for to shiver, then the Arabic for to frighten.
\end{quote}

\subsection*{You'll want to know: \ar{دهشة}}
\textbf{\ar{دهشة}} — \emph{dahsha} — \textquotedblleft{}astonishment\textquotedblright{}.

\begin{quote}
\ar{فتح} \ar{الولد} \ar{الهدية} \ar{بدهشة}.
\end{quote}

\begin{quote}
\emph{fataḥa al-walad al-hadiyya bi-dahsha.}
\end{quote}

\begin{quote}
The boy opened the present in astonishment.
\end{quote}

\begin{cousinweb}
\textbf{\ar{دهشة}} (\emph{dahsha}) is built on the native Arabic root \emph{d-h-sh}, \textquotedblleft{}to be amazed, to be bewildered\textquotedblright{} (Lane). Dahsha is that amazed state.
\end{cousinweb}

\begin{grammarlens}[title={What you've built}]
One more word for describing what you lived and felt.
\end{grammarlens}

\subsection*{Your turn}
\begin{itemize}
\raggedright
  \item \emph{Say it:} \emph{dahsha}
  \item \emph{Say it:} \emph{dahsha}, once more
  \item \emph{Say it:} \emph{fataḥa al-walad al-hadiyya bi-dahsha.} — The boy opened the present in astonishment.
  \item \emph{Recall:} say the Arabic for to shiver, then the Arabic for to frighten, then say \emph{dahsha} again
\end{itemize}

\begin{tcolorbox}[breakable,colback=teal!4,colframe=teal!35!black,title={Before you move on}]
What does \ar{دهشة} mean? (\textquotedblleft{}Astonishment\textquotedblright{}.) Say it once more.
\end{tcolorbox}
\section[\texorpdfstring{\ar{مندهش} (mundahish)}{مندهش (mundahish)}]{\ar{مندهش} (mundahish) — astonished}

\label{lesson:AR-C490-mundahish}



\begin{quote}
Before the new one: say the Arabic for to frighten, then the Arabic for astonishment.
\end{quote}

\subsection*{You'll want to know: \ar{مندهش}}
\textbf{\ar{مندهش}} — \emph{mundahish} — \textquotedblleft{}astonished\textquotedblright{}.

\begin{quote}
\ar{كان} \ar{أبي} \ar{مندهشا} \ar{من} \ar{الخبر}.
\end{quote}

\begin{quote}
\emph{kāna abī mundahishan min al-khabar.}
\end{quote}

\begin{quote}
My father was astonished by the news.
\end{quote}

After \ar{كان} it takes the accusative: \ar{كان} \ar{مندهشا}.

\begin{cousinweb}
\textbf{\ar{مندهش}} (\emph{mundahish}) is built on the native Arabic root \emph{d-h-sh}, \textquotedblleft{}to be amazed, to be bewildered\textquotedblright{} (Lane). Mundahish is the participle of Form VII: one who has been struck with amazement.
\end{cousinweb}

\begin{grammarlens}[title={What you've built}]
One more word for describing what you lived and felt.
\end{grammarlens}

\subsection*{Your turn}
\begin{itemize}
\raggedright
  \item \emph{Say it:} \emph{mundahish}
  \item \emph{Say it:} \emph{mundahish}, once more
  \item \emph{Say it:} \emph{kāna abī mundahishan min al-khabar.} — My father was astonished by the news.
  \item \emph{Recall:} say the Arabic for to frighten, then the Arabic for astonishment, then say \emph{mundahish} again
\end{itemize}

\begin{tcolorbox}[breakable,colback=teal!4,colframe=teal!35!black,title={Before you move on}]
What does \ar{مندهش} mean? (\textquotedblleft{}Astonished\textquotedblright{}.) Say it once more.
\end{tcolorbox}
\section[\texorpdfstring{\ar{اندهش} (indahasha)}{اندهش (indahasha)}]{\ar{اندهش} (indahasha) — to be amazed}

\label{lesson:AR-C490-indahasha}



\begin{quote}
Before the new one: say the Arabic for astonishment, then the Arabic for astonished.
\end{quote}

\subsection*{You'll want to know: \ar{اندهش}}
\textbf{\ar{اندهش}} — \emph{indahasha} — \textquotedblleft{}to be amazed\textquotedblright{}.

\begin{quote}
\ar{اندهشت} \ar{المعلمة} \ar{من} \ar{جواب} \ar{الطالب}.
\end{quote}

\begin{quote}
\emph{indahashat al-muʿallima min jawāb aṭ-ṭālib.}
\end{quote}

\begin{quote}
The teacher was amazed by the student's answer.
\end{quote}

Present: \ar{يندهش}, \emph{yandahish}.

It takes \ar{من} for the cause.

\begin{cousinweb}
\textbf{\ar{اندهش}} (\emph{indahasha}) is built on the native Arabic root \emph{d-h-sh}, \textquotedblleft{}to be amazed, to be bewildered\textquotedblright{} (Lane). It is Form VII, a pattern that often says a thing happened to you.
\end{cousinweb}

\begin{grammarlens}[title={What you've built}]
One more word for describing what you lived and felt.
\end{grammarlens}

\subsection*{Your turn}
\begin{itemize}
\raggedright
  \item \emph{Say it:} \emph{indahasha}
  \item \emph{Say it:} \emph{indahasha}, once more
  \item \emph{Say it:} \emph{indahashat al-muʿallima min jawāb aṭ-ṭālib.} — The teacher was amazed by the student's answer.
  \item \emph{Recall:} say the Arabic for astonishment, then the Arabic for astonished, then say \emph{indahasha} again
\end{itemize}

\begin{tcolorbox}[breakable,colback=teal!4,colframe=teal!35!black,title={Before you move on}]
What does \ar{اندهش} mean? (\textquotedblleft{}To be amazed\textquotedblright{}.) Say it once more.
\end{tcolorbox}
\section[\texorpdfstring{\ar{تفاجأ} (tafājaʾa)}{تفاجأ (tafājaʾa)}]{\ar{تفاجأ} (tafājaʾa) — to be surprised}

\label{lesson:AR-C490-tafajaa}



\begin{quote}
Before the new one: say the Arabic for astonished, then the Arabic for to be amazed.
\end{quote}

\subsection*{You'll want to know: \ar{تفاجأ}}
\textbf{\ar{تفاجأ}} — \emph{tafājaʾa} — \textquotedblleft{}to be surprised\textquotedblright{}.

\begin{quote}
\ar{تفاجأت} \ar{عندما} \ar{رأيت} \ar{صديقي} \ar{في} \ar{المطار}.
\end{quote}

\begin{quote}
\emph{tafājaʾtu ʿindamā raʾaytu ṣadīqī fī al-maṭār.}
\end{quote}

\begin{quote}
I was surprised when I saw my friend at the airport.
\end{quote}

Present: \ar{يتفاجأ}, \emph{yatafājaʾ}.

It takes \ar{بـ} for the cause: \ar{تفاجأ} \ar{بالخبر}.

\begin{cousinweb}
\textbf{\ar{تفاجأ}} (\emph{tafājaʾa}) is built on the native Arabic root \emph{f-j-ʾ}, \textquotedblleft{}to come upon suddenly\textquotedblright{} (Lane). Tafājaʾa is to be come upon suddenly: to be taken by surprise.
\end{cousinweb}

\begin{grammarlens}[title={What you've built}]
One more word for describing what you lived and felt.
\end{grammarlens}

\subsection*{Your turn}
\begin{itemize}
\raggedright
  \item \emph{Say it:} \emph{tafājaʾa}
  \item \emph{Say it:} \emph{tafājaʾa}, once more
  \item \emph{Say it:} \emph{tafājaʾtu ʿindamā raʾaytu ṣadīqī fī al-maṭār.} — I was surprised when I saw my friend at the airport.
  \item \emph{Recall:} say the Arabic for astonished, then the Arabic for to be amazed, then say \emph{tafājaʾa} again
\end{itemize}

\begin{tcolorbox}[breakable,colback=teal!4,colframe=teal!35!black,title={Before you move on}]
What does \ar{تفاجأ} mean? (\textquotedblleft{}To be surprised\textquotedblright{}.) Say it once more.
\end{tcolorbox}
\section[\texorpdfstring{\ar{إعجاب} (iʿjāb)}{إعجاب (iʿjāb)}]{\ar{إعجاب} (iʿjāb) — admiration}

\label{lesson:AR-C490-ijab}



\begin{quote}
Before the new one: say the Arabic for to be amazed, then the Arabic for to be surprised.
\end{quote}

\subsection*{You'll want to know: \ar{إعجاب}}
\textbf{\ar{إعجاب}} — \emph{iʿjāb} — \textquotedblleft{}admiration\textquotedblright{}.

\begin{quote}
\ar{تكلم} \ar{المدير} \ar{عن} \ar{عمل} \ar{الفريق} \ar{بإعجاب}.
\end{quote}

\begin{quote}
\emph{takallama al-mudīr ʿan ʿamal al-farīq bi-iʿjāb.}
\end{quote}

\begin{quote}
The manager spoke about the team's work with admiration.
\end{quote}

It takes \ar{بـ} for what is admired: \ar{إعجاب} \ar{بالفريق}.

\begin{cousinweb}
\textbf{\ar{إعجاب}} (\emph{iʿjāb}) is built on the native Arabic root \emph{ʿ-j-b}, \textquotedblleft{}to wonder, to marvel\textquotedblright{} (Lane). Iʿjāb is the verbal noun of Form IV, to make someone marvel.
\end{cousinweb}

\begin{grammarlens}[title={What you've built}]
One more word for describing what you lived and felt.
\end{grammarlens}

\subsection*{Your turn}
\begin{itemize}
\raggedright
  \item \emph{Say it:} \emph{iʿjāb}
  \item \emph{Say it:} \emph{iʿjāb}, once more
  \item \emph{Say it:} \emph{takallama al-mudīr ʿan ʿamal al-farīq bi-iʿjāb.} — The manager spoke about the team's work with admiration.
  \item \emph{Recall:} say the Arabic for to be amazed, then the Arabic for to be surprised, then say \emph{iʿjāb} again
\end{itemize}

\begin{tcolorbox}[breakable,colback=teal!4,colframe=teal!35!black,title={Before you move on}]
What does \ar{إعجاب} mean? (\textquotedblleft{}Admiration\textquotedblright{}.) Say it once more.
\end{tcolorbox}
